\section{Gradienttimenetelmä}

Gradienttimenetelmä ja sen stokastinen muunnelma ovat ensimmäisen
kertaluvun optimointimenetelmiä, joissa kohdefunktion gradienttia
hyödynnetään parametrivektorin päivittämisessä. Tarkastellaan
minimointitehtävää
\begin{align*}
  \min_{\theta \in \mathbb{R}^d} J(\theta),
\end{align*}
missä \(J : \mathbb{R}^d \to \mathbb{R}\) on jatkuvasti derivoituva
kohdefunktio ja \(\theta \in \mathbb{R}^d\) sen parametrivektori.
Gradienttimenetelmän perusajatus perustuu kohdefunktion ensimmäisen
kertaluvun paikalliseen approksimaatioon. Olkoon
\(\theta_t \in \mathbb{R}^d\) ajanhetken \(t\) parametrivektori ja
\(v \in \mathbb{R}^d\) pieni siirtymävektori. Taylorin lauseen nojalla
saadaan
\begin{align*}
  J(\theta_t + v)
  =
  J(\theta_t)
  +
  \langle \nabla J(\theta_t), v\rangle
  +
  o(\|v\|_2),
  \qquad
  \text{kun } \|v\|_2 \to 0.
\end{align*}

Tästä seuraa, että funktion arvon paikallista muutosta approksimoi
ensisijaisesti lineaarinen termi
\(\langle \nabla J(\theta_t), v\rangle\). Rajoittamalla siirtymän
pituutta ehdolla \(\|v\|_2 \le \varepsilon\) saadaan minimointitehtävä
\begin{align*}
  \min_{\|v\|_2 \le \varepsilon}
  \langle \nabla J(\theta_t), v\rangle.
\end{align*}
Olettamalla, että \(\nabla J(\theta_t) \neq 0\), Cauchy--Schwarzin
epäyhtälöstä seuraa arvio
\begin{align*}
  \langle \nabla J(\theta_t), v\rangle
  \ge
  -
  \|\nabla J(\theta_t)\|_2 \|v\|_2
  \ge
  -
  \varepsilon \|\nabla J(\theta_t)\|_2.
\end{align*}
Yhtäsuuruus toteutuu, kun vektori \(v\) on vastakkaissuuntainen
gradienttiin nähden. Näin ollen lineaarisen approksimaation kannalta
optimaalinen siirtymä on
\begin{align*}
  v^\ast
  =
  -
  \varepsilon
  \frac{\nabla J(\theta_t)}
  {\|\nabla J(\theta_t)\|_2}.
\end{align*}
Tämä motivoi gradienttimenetelmän iteratiivisen päivityssäännön
\begin{align*}
  \theta_{t+1}
  =
  \theta_t - \eta \nabla J(\theta_t),
\end{align*}
missä \(\eta > 0\) on oppimisnopeus (engl. \emph{learning rate}).

Koneoppimisessa kohdefunktio muodostetaan tyypillisesti empiirisenä
riskinä:
\begin{align*}
  J(\theta)
  =
  \frac{1}{N}\sum_{i=1}^N \ell(\theta; z_i),
\end{align*}
missä \(z_1,\dots,z_N\) ovat opetusdatan näytteitä ja
\(\ell(\theta; z_i)\) on havaintoon \(z_i\) liittyvä häviöfunktio.
Tällöin täysi gradientti on
\begin{align*}
  \nabla J(\theta)
  =
  \frac{1}{N}\sum_{i=1}^N
  \nabla \ell(\theta; z_i).
\end{align*}
Kielimallien opetuksessa käytettävät aineistot ovat usein erittäin
suuria, minkä vuoksi täyden gradientin laskeminen jokaisella
iteraatiolla on laskennallisesti kallista. Tämän vuoksi käytetään
stokastista gradienttimenetelmää (engl. \emph{stochastic gradient
descent}, SGD), jossa täysi gradientti
korvataan satunnaisesti poimitun mini-erän perusteella muodostetulla
gradienttiestimaatilla \cite{robbins1951stochastic,bottou1998online,bottou2018optimization}.

Olkoon \(B_t \subset \{1,\dots,N\}\) hetkellä \(t\) muodostettu
satunnainen mini-erä, jonka koko on \(|B_t|=b\). Tällöin gradienttiestimaatti
\(g_t\) määritellään muodossa
\begin{align*}
  g_t
  =
  \frac{1}{b}
  \sum_{i \in B_t}
  \nabla \ell(\theta_t; z_i),
\end{align*}
ja SGD:n päivityssääntö on
\begin{align*}
  \theta_{t+1}
  =
  \theta_t - \eta g_t.
\end{align*}
Jos mini-erä valitaan hetkellä \(t\) tasaisesti kaikkien
\(b\)-alkioisten indeksijoukkojen joukosta, gradienttiestimaatti on
ehdollisen odotusarvon mielessä harhaton \cite{bottou2018optimization}
\begin{align*}
  \mathbb{E}[g_t \mid \theta_t]
  =
  \nabla J(\theta_t).
\end{align*}
Stokastinen gradientti vastaa siten keskimäärin täyttä gradienttia,
vaikka yksittäinen päivitys sisältää satunnaisuudesta johtuvaa kohinaa.
Tämä kohina estää yleensä kohdefunktion arvon monotonisen pienenemisen,
mutta mini-erien käyttö lieventää
laskennallisia haasteita \cite{bottou1998online,goodfellow2016deep}.

Sekä gradienttimenetelmän että SGD:n päivitykset perustuvat
kohdefunktion ensimmäisen kertaluvun paikalliseen approksimaatioon.
Tämä approksimaatio ei kuitenkaan sellaisenaan takaa kohdefunktion
arvon pienenemistä. Tällaisen takuun johtamiseksi tarvitaan gradientin
säännöllisyyttä koskevia lisäoletuksia.

\begin{definition}\label{def:l_smooth}
  Jatkuvasti derivoituva funktio \(J : \mathbb{R}^d \to \mathbb{R}\) on
  \(L\)-sileä, jos sen gradientti \(\nabla J\) on \(L\)-Lipschitz, eli
  kaikilla \(\theta,\theta' \in \mathbb{R}^d\) pätee
  \begin{align*}
    \|\nabla J(\theta) - \nabla J(\theta')\|_2
    \le
    L\|\theta-\theta'\|_2.
  \end{align*}
\end{definition}

Olettaen, että kohdefunktio on \(L\)-sileä, sen muutokselle voidaan
johtaa neliöllinen yläraja.

\begin{lemma}[Neliöllinen yläraja]\label{lemma:quadratic_upper_bound}
  Olkoon \(J:\mathbb{R}^d \to \mathbb{R}\) \(L\)-sileä funktio. Tällöin
  kaikilla \(\theta,h \in \mathbb{R}^d\) pätee
  \begin{align*}
    J(\theta+h)
    \le
    J(\theta)
    +
    \langle \nabla J(\theta), h\rangle
    +
    \frac{L}{2}\|h\|_2^2.
  \end{align*}
\end{lemma}

\begin{proof}
  Analyysin peruslauseesta seuraa, että
  \begin{align*}
    J(\theta+h)-J(\theta)
    =
    \int_0^1
    \langle \nabla J(\theta+th),h\rangle\,dt.
  \end{align*}
  Lisäämällä ja vähentämällä termin
  \(\langle\nabla J(\theta),h\rangle\) integraalin sisällä saadaan
  \begin{align*}
    J(\theta+h)-J(\theta)
    &=
    \int_0^1
    \langle\nabla J(\theta),h\rangle\,dt \\
    &\quad+
    \int_0^1
    \langle\nabla J(\theta+th)-\nabla J(\theta),h\rangle\,dt \\
    &=
    \langle\nabla J(\theta),h\rangle
    +
    \int_0^1
    \langle\nabla J(\theta+th)-\nabla J(\theta),h\rangle\,dt.
  \end{align*}
  Cauchy--Schwarzin epäyhtälön ja \(L\)-Lipschitz-oletuksen perusteella
  integraalissa oleva termi voidaan arvioida ylhäältä:
  \begin{align*}
    \langle\nabla J(\theta+th)-\nabla J(\theta),h\rangle
    &\le
    \|\nabla J(\theta+th)-\nabla J(\theta)\|_2
    \|h\|_2 \\
    &\le
    Lt\|h\|_2^2.
  \end{align*}
  Integroimalla saadaan
  \begin{align*}
    \int_0^1
    \langle\nabla J(\theta+th)-\nabla J(\theta),h\rangle\,dt
    &\le
    \int_0^1 Lt\|h\|_2^2\,dt \\
    &=
    \frac{L}{2}\|h\|_2^2.
  \end{align*}
  Näin ollen
  \begin{align*}
    J(\theta+h)
    \le
    J(\theta)
    +
    \langle\nabla J(\theta),h\rangle
    +
    \frac{L}{2}\|h\|_2^2.
  \end{align*}
\end{proof}

Soveltamalla neliöllistä ylärajaa gradienttimenetelmän päivitykseen
\(h=-\eta\nabla J(\theta)\) saadaan
\begin{align*}
  J(\theta-\eta\nabla J(\theta))
  &\le
  J(\theta)
  -
  \eta\|\nabla J(\theta)\|_2^2
  +
  \frac{L\eta^2}{2}\|\nabla J(\theta)\|_2^2 \\
  &=
  J(\theta)
  -
  \eta\left(1-\frac{L\eta}{2}\right)
  \|\nabla J(\theta)\|_2^2.
\end{align*}
Erityisesti, jos \(0<\eta\le 1/L\), niin
\begin{align*}
  J(\theta-\eta\nabla J(\theta))
  \le
  J(\theta)
  -
  \frac{\eta}{2}\|\nabla J(\theta)\|_2^2.
\end{align*}
Tämä todistaa laskeutumislemman.

\begin{lemma}[Laskeutumislemma]\label{lem:gradient_descent_lemma}
  Olkoon \(J:\mathbb{R}^d\to\mathbb{R}\) \(L\)-sileä funktio ja
  \(0<\eta\le 1/L\). Tällöin gradienttimenetelmän päivityksellä
  \begin{align*}
    \theta_{t+1}
    =
    \theta_t-\eta\nabla J(\theta_t)
  \end{align*}
  pätee
  \begin{align*}
    J(\theta_{t+1})
    \le
    J(\theta_t)
    -
    \frac{\eta}{2}\|\nabla J(\theta_t)\|_2^2.
  \end{align*}
\end{lemma}

Lemma osoittaa, että kohdefunktion arvo pienenee jokaisessa
ei-stationaarisessa pisteessä, kun oppimisnopeus valitaan ehdon
\(0<\eta\le 1/L\) mukaisesti. Käytännössä oppimisnopeuden valinta ei
kuitenkaan ole suoraviivaista, sillä kohdefunktion kaarevuus voi vaihdella
merkittävästi eri parametrisuunnissa. Tällöin yksi kiinteä oppimisnopeus
voi olla joissakin suunnissa liian suuri ja toisissa tarpeettoman pieni.
Tämä motivoi adaptiiviset menetelmät, joissa askelpituutta voidaan mukauttaa optimoinnin aikana parametrikohtaisesti.

\subsection{Häiriöalttius ja optimoinnin haasteet}
\label{subsec:häiriöalttius}

Stokastisen gradienttimenetelmän keskeinen haaste liittyy siihen, että
kohdefunktion lokaali kaarevuus voi vaihdella huomattavasti eri
parametrisuunnissa. Tätä ilmiötä kutsutaan numeeriseksi häiriöalttiudeksi
(engl. \emph{ill-conditioning}). Olkoon kohdefunktio \(J\) \(L\)-sileä ja
kahdesti jatkuvasti derivoituva. Tällöin pisteen \(\theta_t\) ympäristössä
pätee toisen kertaluvun Taylorin approksimaatio
\begin{align*}
  J(\theta_t+\Delta\theta)
  \approx
  J(\theta_t)
  +
  \langle\nabla J(\theta_t),\Delta\theta\rangle
  +
  \frac{1}{2}
  \langle\Delta\theta,\nabla^2J(\theta_t)\Delta\theta\rangle.
\end{align*}
Merkitään Hessin matriisia pisteessä \(\theta_t\) symbolilla
\begin{align*}
  H_t=\nabla^2J(\theta_t).
\end{align*}
Koska \(H_t\) on symmetrinen, sen ominaisarvot ovat reaalisia. Oletetaan
tässä, että lokaali kvadraattinen malli on konveksi, eli
\(H_t\succeq 0\). Tällöin ominaisarvot voidaan järjestää muodossa
\begin{align*}
  \lambda_1\ge\lambda_2\ge\dots\ge\lambda_d\ge 0,
\end{align*}
missä \(\lambda_1\) vastaa suurinta ja \(\lambda_d\) pienintä
paikallista kaarevuutta.

Koska \(J\) on \(L\)-sileä ja kahdesti jatkuvasti derivoituva, Hessin
operaattorinormille pätee
\begin{align*}
  \|H_t\|_2\le L.
\end{align*}
Koska \(H_t\) on symmetrinen, tästä seuraa matriisiepäyhtälö
\begin{align*}
  -LI\preceq H_t\preceq LI,
\end{align*}
missä \(\preceq\) on Löwner-järjestys. Erityisesti lokaalisti konveksisessa
tapauksessa \(H_t\succeq0\) saadaan
\begin{align*}
  0\preceq H_t\preceq LI,
\end{align*}
mistä seuraa
\begin{align*}
  \lambda_1\le L.
\end{align*}
Tämä rajoittaa oppimisnopeuden kokoa, sillä liian suuri askel voi johtaa
epävakaaseen käyttäytymiseen jyrkimmissä suunnissa. Keskeisempi ongelma
on kuitenkin se, että gradienttimenetelmä käyttää samaa oppimisnopeutta
kaikissa parametrisuunnissa. Jos kohdefunktion lokaali kaarevuus vaihtelee
suunnasta toiseen voimakkaasti, oppimisnopeus on valittava suurimman
kaarevuuden ehdoilla. Tällöin menetelmä pysyy vakaana jyrkissä suunnissa,
mutta etenee samalla hitaasti loivissa suunnissa \cite{bertsekas1999nonlinear, nesterov2004introductory}.

Tätä voidaan havainnollistaa tarkastelemalla funktion käyttäytymistä
stationaarisen pisteen \(\theta^\ast\) ympärillä, missä
\(\nabla J(\theta^\ast)=0\). Jos \(J\) on kahdesti jatkuvasti derivoituva,
gradienttia voidaan pisteen \(\theta^\ast\) läheisyydessä approksimoida
lineaarisesti:
\begin{align*}
  \nabla J(\theta_t)
  \approx
  \nabla^2J(\theta^\ast)(\theta_t-\theta^\ast).
\end{align*}
Merkitään \(H^\ast=\nabla^2J(\theta^\ast)\) ja
\(e_t=\theta_t-\theta^\ast\). Tällöin gradienttimenetelmän päivitys saa
lokaalisti approksimaation
\begin{align*}
  e_{t+1}
  \approx
  (I-\eta H^\ast)e_t.
\end{align*}

Koska \(H^\ast\) on symmetrinen, se voidaan diagonalisoida ortonormaalissa
ominaisvektorikannassa. Oletetaan lisäksi, että \(\theta^\ast\) on lokaali
minimi, jolloin \(H^\ast\succeq0\). Jos ominaisarvot järjestetään muotoon
\begin{align*}
  \lambda_1\ge\lambda_2\ge\dots\ge\lambda_d\ge0,
\end{align*}
niin virheen komponentit toteuttavat approksimatiivisesti relaation
\begin{align*}
  e_{t+1}^{(k)}
  \approx
  (1-\eta\lambda_k)e_t^{(k)},
  \qquad
  k=1,\dots,d.
\end{align*}
Sama oppimisnopeus \(\eta\) vaikuttaa siis eri suuntiin eri tavoin.
Suurta ominaisarvoa vastaavissa suunnissa tekijä
\(1-\eta\lambda_k\) muuttuu nopeasti, joten vakaus edellyttää pientä
oppimisnopeutta. Jos \(\eta\) valitaan lähelle arvoa
\(1/\lambda_1\), jyrkimmässä suunnassa pätee
\(1-\eta\lambda_1\approx0\), jolloin eteneminen on nopeaa. Loivimmassa
suunnassa pätee tällöin puolestaan
\(1-\eta\lambda_d\approx1\), mikä merkitsee hidasta supistumista.
Menetelmä voi siten olla jyrkissä suunnissa lähes optimaalinen mutta
loivissa suunnissa samalla hyvin hidas.

Näin gradienttimenetelmä voi joutua tilanteeseen, jossa oppimisnopeus on
jyrkkien suuntien kannalta lähes suurin sallittu, mutta loivissa suunnissa
eteneminen on silti hidasta. Juuri tämä kaarevuuden epätasaisuus tekee
menetelmästä tehottoman huonosti ehdollistuneissa ongelmissa. Ilmiötä
voidaan kuvata Hessin matriisin ehtoluvulla
\begin{align*}
  \kappa
  =
  \frac{\lambda_1}{\lambda_d},
\end{align*}
kun \(\lambda_d>0\). Positiivisesti definiitillä kvadraattisella
kohdefunktiolla gradienttimenetelmän suppenemisnopeus riippuu
oleellisesti ehtoluvusta. Suuri ehtoluku tarkoittaa, että
kaarevuus vaihtelee huomattavasti eri suunnissa, jolloin yksi globaali
oppimisnopeus soveltuu väistämättä huonosti ainakin osaan suunnista
\cite{nocedal2006numerical,saad2003iterative}.

Vaikka analyysi perustuu lokaaliin toisen kertaluvun approksimaatioon, sen johtopäätös on yleisempi. Jos kohdefunktion paikallinen kaarevuus vaihtelee voimakkaasti parametriavaruuden eri suunnissa, yksi kaikissa suunnissa käytettävä oppimisnopeus ei pysty mukautumaan tehokkaasti näihin eroihin. Tällöin oppimisnopeus on valittava jyrkimpien suuntien ehdoilla, mikä hidastaa etenemistä loivemmissa suunnissa \cite{nocedal2006numerical, bertsekas1999nonlinear}. Tämä motivoi menetelmiä, joissa gradienttipäivityksiä skaalataan tai mukautetaan parametrikohtaisesti esimerkiksi gradienttihistorian perusteella.

\subsection{Adaptiiviset optimointimenetelmät}

Edellisessä alaluvussa todettiin, että stokastinen gradienttimenetelmä
käyttää samaa globaalia oppimisnopeutta kaikissa parametrisuunnissa.
Tämä voi olla tehotonta erityisesti silloin, kun kohdefunktion lokaali
kaarevuus tai gradienttien mittakaava vaihtelee voimakkaasti parametrien
välillä. Tällöin vakaa oppimisnopeus on valittava jyrkkien suuntien
ehdoilla, mikä hidastaa etenemistä loivemmissa suunnissa. Adaptiivisten
optimointimenetelmien keskeinen idea on lieventää tätä ongelmaa
skaalaamalla päivitystä parametrikohtaisesti gradienttihistorian
perusteella.

Yleinen adaptiivinen menetelmä voidaan kirjoittaa muodossa
\begin{align*}
  \theta_{t+1}
  =
  \theta_t-\eta_tD_tg_t,
\end{align*}
missä \(g_t\) on hetkellä \(t\) laskettu stokastinen gradientti,
\(\eta_t>0\) on globaali oppimisnopeus ja
\(D_t\in\mathbb{R}^{d\times d}\) on tavallisesti diagonaalinen,
positiivinen skaalausmatriisi, joka riippuu aiemmista gradienteista.
Tällöin kunkin parametrin efektiivinen askelpituus määräytyy sekä
globaalin oppimisnopeuden että kyseisen parametrin gradienttihistoriaan
perustuvan skaalauksen perusteella. Intuitiivisesti komponentit, joiden
gradientit ovat usein suuria, voidaan päivittää varovaisemmin, kun taas
pienempiä tai harvinaisempia gradientteja saaville komponenteille voidaan
sallia suhteellisesti suurempia askelia.

Adaptiivisten menetelmien kehitys syväoppimisessa on rakentunut
erityisesti AdaGradin, RMSPropin ja Adamin ympärille. AdaGrad tuo esiin
parametrikohtaisen oppimisnopeuden perusidean, RMSProp korjaa AdaGradin
käytännöllistä heikkoutta ja Adam yhdistää adaptiivisen skaalaamisen sekä
momentumin. Modernissa neuroverkkojen ja erityisesti
transformer-arkkitehtuurien opetuksessa myös AdamW on saavuttanut
vakiintuneen aseman
\cite{duchi2011adaptive,tieleman2012lecture,kingma2015adam,
loshchilov2019decoupled, vaswani2017attention}.

\paragraph{AdaGrad.}

AdaGrad (\emph{adaptive gradient}) oli ensimmäisiä laajasti tunnettuja
menetelmiä, joissa oppimisnopeus mukautetaan parametrikohtaisesti
gradienttihistorian perusteella \cite{duchi2011adaptive}. Menetelmässä
ylläpidetään neliöityjen gradienttien kumulatiivista summaa
\begin{align*}
  s_t
  =
  \sum_{\tau=1}^t g_\tau^2,
\end{align*}
missä potenssi- ja neliöjuurioperaatiot sekä jakolasku tulkitaan
komponenttikohtaisesti. Päivitys kirjoitetaan muodossa
\begin{align*}
  \theta_{t+1}
  =
  \theta_t
  -
  \eta
  \frac{g_t}{\sqrt{s_t}+\varepsilon},
\end{align*}
missä \(\varepsilon>0\) on pieni vakio numeerisen stabiiliuden
varmistamiseksi.

AdaGradin keskeinen ominaisuus on, että parametreille, joiden gradientit
ovat toistuvasti suuria, kertyy nopeasti suuri nimittäjä, jolloin niiden
efektiivinen oppimisnopeus pienenee. Vastaavasti parametrit, joiden
gradientit ovat harvinaisia tai pieniä, voivat saada suhteellisesti
suurempia päivityksiä. Tämän vuoksi menetelmä soveltuu erityisesti
tilanteisiin, joissa gradientit ovat harvoja. AdaGradin keskeinen
heikkous on kuitenkin se, että kertymä \(s_t\) kasvaa monotonisesti.
Tämän seurauksena askelpituudet voivat ajan myötä pienentyä liikaa ja
optimointi hidastua olennaisesti.

\paragraph{RMSProp.}

RMSProp (\emph{root mean square propagation}) kehitettiin käytännölliseksi
korjaukseksi AdaGradin liian nopeasti pieneneville askelpituuksille.
Menetelmää ei alun perin esitelty vertaisarvioidussa tieteellisessä
artikkelissa, vaan se tuotiin julkisuuteen Geoffrey Hintonin pitämällä
Coursera-verkkokurssilla \cite{tieleman2012lecture}. RMSPropin
toimintaperiaate on sittemmin dokumentoitu laajasti alan perusteoksissa
\cite{goodfellow2016deep}.

Sen sijaan, että neliöityjä gradientteja
summattaisiin koko optimoinnin ajalta, RMSProp käyttää eksponentiaalisesti
painotettua liukuvaa keskiarvoa
\begin{align*}
  v_t
  =
  \rho v_{t-1}+(1-\rho)g_t^2,
\end{align*}
missä \(\rho\in(0,1)\) on muistiparametri. Päivityssääntö on
\begin{align*}
  \theta_{t+1}
  =
  \theta_t
  -
  \eta
  \frac{g_t}{\sqrt{v_t}+\varepsilon}.
\end{align*}

RMSProp säilyttää AdaGradin perusintuition parametrikohtaisesta
skaalaamisesta, mutta käyttää pitkän aikavälin kertymän sijasta
paikallisempaa estimaattia gradienttien mittakaavasta. Tämä estää
oppimisnopeuden systemaattisen pienentymisen ja tekee menetelmästä
käytännössä käyttökelpoisemman syvien neuroverkkojen opetuksessa \cite{goodfellow2016deep}.

\paragraph{Adam.}

Adam (\emph{adaptive moment estimation}) yhdistää kaksi ideaa: gradientin
ensimmäisen momentin eli suunnatun liukuvan keskiarvon sekä neliöityjen
gradienttien toisen momentin eli liukuvan keskiarvon
\cite{kingma2015adam}. Menetelmässä määritellään rekursiot
\begin{align*}
  m_t
  &=
  \beta_1m_{t-1}+(1-\beta_1)g_t, \\
  v_t
  &=
  \beta_2v_{t-1}+(1-\beta_2)g_t^2,
\end{align*}
missä \(\beta_1,\beta_2\in(0,1)\) ovat muistiparametreja. Koska \(m_t\)
ja \(v_t\) alustetaan tavallisesti nollaksi, ne ovat alkuvaiheessa
harhaisia kohti nollaa. Tämän vuoksi käytetään harhakorjattuja estimaatteja
\begin{align*}
  \hat{m}_t
  &=
  \frac{m_t}{1-\beta_1^t}, \\
  \hat{v}_t
  &=
  \frac{v_t}{1-\beta_2^t}.
\end{align*}
Päivitys saa tällöin muodon
\begin{align*}
  \theta_{t+1}
  =
  \theta_t
  -
  \eta
  \frac{\hat{m}_t}{\sqrt{\hat{v}_t}+\varepsilon}.
\end{align*}

Adamia voidaan tulkita menetelmäksi, jossa \(m_t\) vakauttaa
päivityssuuntaa mo-mentum-tyyppisesti, kun taas \(v_t\) skaalaa
askelpituutta parametrikohtaisesti gradienttien havaittuun mittakaavaan
nähden. Tämä tekee menetelmästä käytännössä usein tehokkaan ja
suhteellisen robustin oppimisnopeuden valinnalle. Näistä syistä Adamista
on tullut yksi käytetyimmistä optimointimenetelmistä syväoppimisessa.

Adam ei kuitenkaan ole ongelmaton. On havaittu, että adaptiiviset
menetelmät eivät kaikissa tilanteissa tuota yhtä hyvää yleistymistä kuin
huolellisesti viritetty SGD momentumilla \cite{wilson2017marginal}. Lisäksi Adamin alkuperäinen algoritmi voi tietyillä gradienttisekvensseillä epäonnistua konvergenssissa, minkä vuoksi sen konvergenssianalyysiä on myöhemmin tarkennettu esimerkiksi AMSGrad-menetelmän avulla \cite{reddi2018convergence}.

Adaptiivisuutta ei siten
tule tulkita yleispätevänä ratkaisuna kaikkiin optimoinnin haasteisiin,
vaan käytännöllisenä keinona muokata päivityksen geometriaa ja vähentää
oppimisnopeuden herkkyyttä.

\paragraph{AdamW.}

AdamW on Adamin muunnelma, jossa painojen kutistaminen
(\emph{weight decay}) erotetaan adaptiivisesta gradienttipäivityksestä
\cite{loshchilov2019decoupled}. Tavallisessa \(L_2\)-regulari-soinnin ja
Adamin yhdistelmässä regularisaatiotermi sekoittuu adaptiiviseen
skaalaukseen tavalla, joka ei vastaa puhdasta painojen kutistamista.
AdamW:ssä parametripäivitys kirjoitetaan tyypillisesti muodossa
\begin{align*}
  \theta_{t+1}
  =
  \theta_t
  -
  \eta
  \frac{\hat{m}_t}{\sqrt{\hat{v}_t}+\varepsilon}
  -
  \eta\lambda\theta_t,
\end{align*}
missä \(\lambda\ge0\) säätelee painojen kutistamisen voimakkuutta.

Tämä erotus on osoittautunut käytännössä tärkeäksi erityisesti syvien
neuroverkkojen opetuksessa \cite{loshchilov2019decoupled}. AdamW:tä
pidetäänkin usein tarkoituksenmukaisempana valintana kuin alkuperäistä
Adamia silloin, kun mallin opetuksessa käytetään eksplisiittistä painojen
kutistamista. Transformer-arkkitehtuurien opetuksessa Adam-tyyppiset
adaptiiviset menetelmät ovat yleisesti käytettyjä
\cite{vaswani2017attention}.

Lopuksi on syytä huomata, että hyvin suurissa kielimalleissa myös
optimointimenetelmän muistikustannus muodostuu olennaiseksi. Adam
ylläpitää jokaiselle parametrille ensimmäisen ja toisen momentin
estimaatit, minkä vuoksi optimointitila voi kuluttaa huomattavan määrän
muistia. Tämän vuoksi on kehitetty muistitehokkaampia adaptiivisia
menetelmiä, kuten Adafactor, joissa pyritään säilyttämään adaptiivisen
skaalaamisen keskeiset edut pienemmällä muistinkulutuksella
\cite{shazeer2018adafactor}.

Tässä tutkielmassa tarkastelun painopiste rajataan kuitenkin
AdaGrad-, RMS-Prop- ja Adam-tyyppisiin menetelmiin, sillä ne muodostavat
historiallisesti ja käsitteellisesti keskeisen perustan myöhemmille
adaptiivisille optimointialgoritmeille.
